% ID: SEL-D0049-OA-E01
% ESTADO: cotejada. Segunda pasada visual y borrador compilado revisados.
% Fuente: Selectividad/sel_2013_matematicas_aplicadas/junio/reserva b/EXAMEN 1-ANDALUCÍA.pdf
% SHA-256: ec6fd00d0daee547aa6128f3f6e1099278ee551f7f80917b331a1772764317b3
% Página: 1; opción A; ejercicio original 1.

\textbf{(2.5 puntos)} Un fabricante elabora dos tipos de anillos a base de oro y plata. Cada anillo del primer tipo precisa 4 g de oro y 2 de plata, mientras que cada uno del segundo necesita 3 g de oro y 1 de plata. Sabiendo que dispone de 48 g de oro y 20 de plata y que los precios de venta de cada tipo de anillo son 150 euros el primero y 100 euros el segundo, \textquestiondown{}cuántos anillos de cada tipo tendría que producir para obtener los ingresos máximos? \textquestiondown{}A cuánto ascenderían estos ingresos?
